\section{Euler's \texorpdfstring{$\phi$}{phi} Function}
\begin{definition}[Euler's Totient Function, $\phi$]
    Let $m \in \mathbb{N}$ such that $m \ge 1$. The **Euler's totient function**, denoted by $\phi(m)$, is the number of elements in the multiplicative group of integers modulo $m$, which is $\left(\mathbb{Z}/m\mathbb{Z}\right)^*$.
    
    In other words, $\phi(m)$ is the number of integers $k$ such that $1 \le k \le m$ so that $\gcd(k, m) = 1$.
\end{definition}

\begin{center}
    \vspace{10pt}
    \line(1,0){450}
    \vspace{25pt}
\end{center}

\begin{definition}[Remainder Map]
    Let $m, n \in \mathbb{N}$ be integers. The **remainder map** $r$ is a function:
    $$r: \mathbb{Z}/(mn)\mathbb{Z} \longrightarrow (\mathbb{Z}/m\mathbb{Z}) \times (\mathbb{Z}/n\mathbb{Z})$$
    
    It is defined $\forall a \in \{0, \dots, mn - 1\}$ as:
    $$r([a]_{mn}) := ([a]_m, [a]_n)$$
    where $[a]_k$ denotes the residue class of $a$ modulo $k$.
\end{definition}

\begin{center}
    \vspace{10pt}
    \line(1,0){450}
    \vspace{25pt}
\end{center}

\bigskip

\begin{proposition}[Remainder Map is Well-Defined]
    The remainder map $r$ is well-defined. That is, the result of the map is independent of the choice of the representative of the residue class:
    $$[a]_{mn} = [b]_{mn} \implies r([a]_{mn}) = r([b]_{mn})$$
    which means $([a]_m, [a]_n) = ([b]_m, [b]_n)$.
\end{proposition}

\begin{proof}
    We assume the antecedent: $[a]_{mn} = [b]_{mn}$.
    
    By the definition of equality of residue classes, this means $mn$ divides the difference $(b - a)$:
    $$[a]_{mn} = [b]_{mn} \implies mn \mid (b - a)$$
    
    Since $mn \mid (b - a)$, it must be that both $m$ and $n$ individually divide the difference $(b - a)$:
    $$mn \mid (b - a) \implies \begin{cases} m \mid (b - a) \\ n \mid (b - a) \end{cases}$$
    
    By the definition of congruence and residue classes, this translates back into the desired equalities:
    $$\begin{cases} m \mid (b - a) \implies [a]_m = [b]_m \\ n \mid (b - a) \implies [a]_n = [b]_n \end{cases}$$
    
    Since both components are equal, the ordered pairs are equal:
    $$([a]_m, [a]_n) = ([b]_m, [b]_n)$$
    
    Thus, $r([a]_{mn}) = r([b]_{mn})$, and the map is well-defined.
\end{proof}

\begin{center}
    \vspace{10pt}
    \line(1,0){450}
    \vspace{25pt}
\end{center}


\begin{theorem}[Bijectivity of the Remainder Map with Coprime Factors]
    Let $m, n \in \mathbb{N}$ be integers such that $m, n \ge 1$, and assume $m$ and $n$ are **pairwise coprime** (i.e., $\gcd(m, n) = 1$).
    
    The remainder map $r$:
    $$r: \mathbb{Z}/(mn)\mathbb{Z} \longrightarrow (\mathbb{Z}/m\mathbb{Z}) \times (\mathbb{Z}/n\mathbb{Z})$$
    is a **bijection** (one-to-one and onto).
\end{theorem}

\begin{proof}
    The proof establishes bijectivity by demonstrating both **Surjectivity** (the map is onto) and **Injectivity** (the map is one-to-one).

    \subsection*{Surjectivity (Onto)}
    Let $([a]_m, [b]_n) \in (\mathbb{Z}/m\mathbb{Z}) \times (\mathbb{Z}/n\mathbb{Z})$ be an arbitrary element in the codomain.
    
    We seek to show $\exists x \in \mathbb{Z}$ so that $r([x]_{mn}) = ([a]_m, [b]_n)$.
    
    This is equivalent to finding an integer $x$ that simultaneously solves the system of congruences:
    $$\begin{cases} x \equiv a \pmod{m} \\ x \equiv b \pmod{n} \end{cases}$$
    
    Since we are given that $\gcd(m, n) = 1$, the **Chinese Remainder Theorem** guarantees that a solution $x$ always exists for any pair of residues $(a, b)$.
    
    Therefore, $\forall$ element in the codomain, there is a corresponding pre-image in the domain, proving that the map $r$ is **surjective**.
    
    \subsection*{Injectivity (One-to-One)}
    We can establish injectivity by counting the elements (since the sets are finite):
    
    \begin{enumerate}
        \item **Count the Domain:** The number of elements in the domain is the modulus:
        $$|\mathbb{Z}/(mn)\mathbb{Z}| = mn$$
        
        \item **Count the Codomain:** The number of elements in the codomain (a Cartesian product) is the product of the sizes of the factors:
        $$|(\mathbb{Z}/m\mathbb{Z}) \times (\mathbb{Z}/n\mathbb{Z})| = |\mathbb{Z}/m\mathbb{Z}| \cdot |\mathbb{Z}/n\mathbb{Z}| = m \cdot n = mn$$
        
        \item **Conclusion:** Since the domain and codomain are finite sets and have the **same number of elements** ($mn = mn$), and we have already proven the map $r$ is **surjective**, it must automatically be **injective** as well.
        
        $$\text{Surjectivity} \land (|\text{Domain}| = |\text{Codomain}|) \implies \text{Injectivity}$$
    \end{enumerate}
    
    Since $r$ is both surjective and injective, it is a **bijection**.
\end{proof}

\begin{center}
    \vspace{10pt}
    \line(1,0){450}
    \vspace{25pt}
\end{center}


\begin{theorem}[Totient of a Prime Number]
\label{thm:totient_prime}
The Totient of a Prime Number: If $m \in \mathbb{N}$ so that $m$ is prime, then $\varphi(m) = m - 1$.
\end{theorem}

\begin{center}
    \vspace{10pt}
    \line(1,0){450}
    \vspace{25pt}
\end{center}

\begin{theorem}[Primality Test via the Totient Function]
\label{thm:primality_test}
A Primality Test via the Totient Function: $\forall m \in \mathbb{N} : m \geq 2$, $\varphi(m) = m - 1 \iff m$ is a prime number.
\end{theorem}
\begin{proof}
    \leavevmode
\begin{itemize}
    \item[($\Rightarrow$)] \textbf{(``Only if'' direction)}: We prove the contrapositive. Assume $m$ is composite (not prime). $\implies \exists p \in \mathbb{N} : 1 < p < m \text{ and } p \mid m$.
    Since $p \mid m$, $\gcd(p, m) = p > 1$. Thus, $p$ is not coprime to $m$. Since $p \in \{1, \dots, m-1\}$ is excluded from the count, the total count $\varphi(m)$ must be strictly less than $m-1$. More precisely, $\varphi(m) \leq m - 2$. This shows that the equality $\varphi(m) = m - 1$ holds $\implies m$ is prime.
    \item[($\Leftarrow$)] \textbf{(``If'' direction)}: Assume $m$ is prime. Then $\forall a \in \{1, \dots, m-1\}$, $\gcd(a, m) = 1$. Since there are $m-1$ such integers, $\varphi(m) = m - 1$. \qedhere
\end{itemize}
\end{proof}

\begin{center}
    \vspace{10pt}
    \line(1,0){450}
    \vspace{25pt}
\end{center}

\begin{proposition}[Totient of a Prime Power]
\label{prop:prime_power_formula}
$\forall p$ prime and $\forall \alpha \in \mathbb{N} : \alpha \geq 1$:
$$ \varphi(p^\alpha) = p^\alpha - p^{\alpha-1} $$
\end{proposition}
\begin{proof}
The totient function $\varphi(p^\alpha)$ counts the number of integers $a \in \{1, \dots, p^\alpha\}$ so that $\gcd(a, p^\alpha) = 1$.
The integers that are \textbf{not} coprime to $p^\alpha$ are those $a$ for which $p \mid a$. We count these integers (the complement).
These non-coprime integers are the multiples of $p$ in the range $[1, p^\alpha]$, i.e., of the form $k \cdot p$ where $k \in \mathbb{N}$.
The condition $1 \leq k \cdot p \leq p^\alpha$ implies $1 \leq k \leq p^{\alpha-1}$.
Thus, there are exactly $p^{\alpha-1}$ multiples of $p$ in the specified range.
Subtracting this count from the total number of integers, $p^\alpha$, yields:
$$ \varphi(p^\alpha) = p^\alpha - p^{\alpha-1} = p^{\alpha-1} (p - 1) $$ \qedhere
\end{proof}

\begin{center}
    \vspace{10pt}
    \line(1,0){450}
    \vspace{25pt}
\end{center}

\begin{theorem}[Multiplicativity of Euler's Totient Function]
\label{thm:multiplicativity}
Multiplicativity of Euler's Totient Function: If $m, n \in \mathbb{N}$ are coprime integers ($\gcd(m, n) = 1$), then $\varphi(m n) = \varphi(m) \varphi(n)$.
\end{theorem}